\section{Test Harness and Hardware-in-the-Loop Demonstration}
\label{sec:hil}

This chapter describes the hardware-in-the-loop (HIL) test
infrastructure built around the firmware: the multi-tier
architecture (web app, ESP32 bridge, SAMD21 board), the demo pages
exposed to the operator, the twelve preset scenarios used to
exercise the state machine, the MPPT convergence demo, the
real-mainboard bring-up, and the bench-test programme that
produces the physical-measurement evidence backing each firmware
claim.

\subsection{Why hardware-in-the-loop}

A pure-software simulator can prove that an algorithm produces
the right answer for a clean set of inputs. It cannot prove that
the algorithm produces the right answer \textit{on the actual
chip}, with real ADC noise, real interrupt latency, real
peripheral configuration, and real wiring. Conversely, a pure
hardware test --- a board with a solar panel, a battery, and a
load --- is expensive to set up, hard to script, and gives slow
feedback when an algorithm needs adjustment.

The HIL approach takes a middle path. The firmware runs unchanged
on the real chip on the real board. The sensor environment ---
the solar panel and the battery --- is replaced by a software
model running on an external companion microcontroller. The
companion exchanges raw ADC values with the SAMD21 over a serial
link, closing the loop in real time. Anything the firmware does
on real hardware (peripheral configuration, interrupt handling,
timing) is exercised. Anything the firmware does with the
environment (algorithm decisions, state transitions) is exercised
against scriptable, repeatable stimuli.

\subsection{The architecture of the harness}

The harness has three tiers:

\begin{description}
\item[The laptop.] A Python web application serves three local
pages (\texttt{/mppt}, \texttt{/state}, \texttt{/manual}) at
\texttt{http://127.0.0.1:8000}. Each page renders an HTML/CSS/JS
front-end and posts text commands to a small set of
\texttt{/api/...} endpoints. The application owns a single
USB-serial connection.

\item[The ESP32 bridge.] An Arduino sketch on an ESP32
dev-board acts as the OBC stand-in. It translates the plain-text
commands forwarded from the laptop into binary CHIPS frames, sends
them to the SAMD21 over a 3.3\,V UART link, parses the replies,
and prints them back to the laptop in a regex-friendly format.

\item[The SAMD21.] The real EPS firmware. It receives CHIPS frames
on SERCOM0, dispatches commands, runs the selected mode, applies
outputs, and streams status replies back to the bridge.
\end{description}

\begin{figure}[ht]
\centering
\fbox{\parbox{0.85\textwidth}{\centering [Figure placeholder]\\
Three-tier HIL architecture: browser $\rightarrow$ Python web
app (\texttt{run\_local\_web\_app.py}) $\rightarrow$ USB serial
(\texttt{COM3}) $\rightarrow$ ESP32 bridge $\rightarrow$ 3.3\,V
UART $\rightarrow$ SAMD21 mainboard. CHIPS frames flow in both
directions on the 3.3\,V link; plain text on the USB link.}}
\caption{The three-tier hardware-in-the-loop architecture.}
\label{fig:hil-architecture}
\end{figure}

The advantage of routing everything through one port and one
Python process is that no second connection has to be opened
when a new page is added: a new page becomes a new actions module
and a new \texttt{static/<page>/} folder; the existing serial
link is reused. The HTML pages share a single live-telemetry
sidebar that updates from the streamed snapshots.

\subsection{The three demo pages}

\subsubsection{MPPT convergence}

The \texttt{/mppt} page lets the operator configure a simulated
solar-panel I--V curve (open-circuit voltage, short-circuit
current, voltage at the maximum power point), launch the MPPT
demo, and watch the algorithm converge in real time. The ESP32
runs the physics --- the simulated panel and the buck converter
model --- and the SAMD21 runs the algorithm. The page plots
panel voltage, panel current, panel power, and commanded duty
cycle on a shared time axis, marking the theoretical maximum-
power point as a horizontal line. Convergence is visually
obvious to a reviewer who is not familiar with the algorithm.

\subsubsection{State-transition scenarios}

The \texttt{/state} page exposes twelve preset scenarios that
each inject a steady set of sensor values and assert which
state the firmware should reach. The scenarios cover the four
normal PCU modes (\texttt{MPPT\_CHARGE},
\texttt{CV\_FLOAT}, \texttt{SA\_LOAD\_FOLLOW},
\texttt{BATTERY\_DISCHARGE}), the three safe-mode entry triggers
(battery low, temperature out of range, heartbeat lost), and the
four safe-mode sub-state load masks. Running a scenario sends
the sensor values, waits a few seconds for steady state, reads
back the firmware's snapshot, and compares the observed mode and
output mask to the expected values. A green dot means the
scenario passed; a red dot opens a side-by-side table of
expected versus observed values.

A run-all button executes the suite in sequence and reports a
single pass/fail count. The current implementation reaches
12 of 12 passes against the pure-dispatcher firmware running on
the real board.

\subsubsection{Manual control}

The \texttt{/manual} page lets the operator drive each of the
four user-facing board outputs independently: the PWM duty
cycle, the photovoltaic (PV) eFuse enable, the battery (BAT)
eFuse enable, and the green status LED. The page enforces an
``Off first'' convention before allowing entry into manual
mode and disarms the PWM on every disconnect. Manual mode is
primarily used during hardware bring-up when a specific output
needs to be observed in isolation.

A fourth page, \texttt{/comm} (communication and status), is
listed in the project plan but is not yet implemented.

\subsection{Real-mainboard bring-up}

The CHESS EPS firmware was first developed against the
Curiosity Nano DM320119 dev board (SAMD21G17D) and then ported
to the real PCU testing board V4.1 (SAMD21J17D-MUT). The two
chips are the same silicon die in different packages, so the
register-level code is unchanged; only the pin assignments,
the chip define, the startup file, and the linker script differ
between the two builds.

The mainboard bring-up established:

\begin{itemize}
\item the programming path: Atmel-ICE through a Tag-Connect
TC2050-IDC-NL footprint at J1, with SWD clock at 125\,kHz;
\item the programming software path: Microchip Studio Device
Programming (the custom-cable adapter does not connect reset,
so reset-dependent flows could not be used);
\item the CHIPS link: the ESP32 bridge wired to J3 (\texttt{TX}
to PA11, \texttt{RX} to PA10, common GND) exchanges valid
frames at 115200\,baud;
\item the full firmware (\texttt{satellite\_firmware.hex}) runs
on the real chip; telemetry, state-machine scenarios, MPPT
demo, and PWM safety gating all behave as on the dev board;
\item the buck PWM signals on PA12/PA13 were observed on an
oscilloscope, with a measured dead-time gap of approximately
100\,ns at the first checked edge.
\end{itemize}

The second PWM edge and a full duty-cycle sweep are part of the
planned bench programme (Section~\ref{sec:hil}.\ref{sec:bench-test}).

\subsection{The bench-test programme}
\label{sec:bench-test}

The HIL harness exercises algorithm logic. To complete the
verification chain the firmware also needs to be exercised
against \textit{physical electrical stimuli}: known voltages
applied to ADC inputs, known currents through the shunt
resistors, a real buck converter producing real V$_{\text{out}}$,
and a real fake-solar source whose maximum power point can be
computed by hand.

A structured nine-test programme has been planned and is
documented in detail in \texttt{src-pds/how\_to\_test.md}.
Table~\ref{tab:bench-tests} summarises each test, its
dependencies, and the evidence it produces for the report.

\begin{table}[ht]
\centering
\caption{The nine bench tests and what each produces.}
\label{tab:bench-tests}
\small
\begin{tabular}{@{}p{0.05\linewidth}p{0.4\linewidth}p{0.5\linewidth}@{}}
\toprule
\textbf{ID} & \textbf{Test} & \textbf{Evidence produced} \\
\midrule
A & PWM waveform sign-off & Scope shots of both edges at three
duty cycles; dead-time table; frequency measurement. \\
B & ADC calibration vs known voltage & Calibration table per
sensor net; observed vs reference ADC reference voltage. \\
C & eFuse switching on real board & Multimeter pre/post
readings; PGOOD readback table; verified GPIO actuation. \\
D & State machine on real sensors & Mode-transition table with
hardware-fed inputs; cross-check against the 12-scenario suite. \\
E & MPPT against a fake solar source ($\star$ headline) &
Convergence plot from a known Thevenin source; comparison to the
hand-calculated maximum power point. \\
F & Buck transfer curve & Duty vs $V_{\text{out}}/V_{\text{in}}$
plot showing the expected straight-line relationship. \\
G & INA226 sensor verification & Bus-scan output; ID
register reads; multimeter cross-check at two current points. \\
H & CHIPS communication stress & Total-frames count after a
ten-minute stream; bad-CRC and duplicate-sequence behaviour. \\
I & Watchdog timer & Uptime trace around a deliberate stall;
verification of reset and recovery. \\
\bottomrule
\end{tabular}
\end{table}

The plan also captures the dependencies between the tests (Test G
requires the I\textsuperscript{2}C address fix; Test E requires
the fake-solar power resistors; Tests B, D and the MPPT-on-real-
sensors workflow require a host-reachable
\texttt{set\_sensor\_source} command). Each test prescribes the
exact silkscreen designators, test points, and PCB coordinates
for every probe location, derived from a primary-source read of
the public PCB schematic and BOM (Chapter~\ref{sec:results}
discusses the research process). The user running the test does
not need to read the schematic itself.

When the bench programme is executed, the resulting tables,
plots, and scope shots will populate the figures in
Chapter~\ref{sec:results}.
